\section*{Anexo 4-D — Matriz de los doce módulos}
\phantomsection\label{fuente-05-anexo-4-1-d-matriz-de-los-doce-modulos}
\phantomsection
\label{anx:D}
\addcontentsline{toc}{section}{Anexo 4-D — Matriz de los doce módulos}
La matriz asigna a cada módulo su requisito funcional, su responsabilidad, su interfaz principal, su actor y la etapa en que entra en producción.

\begin{tablalafrox}{Módulos, responsabilidades e interlocutores}{tab:modulos-funcionales}{F{0.18} F{0.23} F{0.24} F{0.23} Y}{\cab{Módulo / RF} & \cab{Responsabilidad} & \cab{Interfaz principal} & \cab{Actor principal} & \cab{Etapa}}
M1 Recepción / RF-01 & Lote y recepción GS1. & Evento a M2, ACL con ERP. & Jefa de Bodega, Proveedor. & 1 \\
M2 Inventario / RF-02 y RF-10 & Movimientos por sitio, reserva central, FEFO y sugerencia de reposición de compras por SKU y proveedor con promociones comprometidas. & Consulta de M3/M5, sugerencia aprobada al ERP por la ACL de INT-06. & Jefa de Bodega, función de abastecimiento dentro de los perfiles autorizados. & 1 \\
M3 Preventa / RF-03 & Pedido y precio pactado, autoatención desde el Portal de Clientes. & Reserva coordinada en M2, sincronización. & Preventista, clientes autorizados. & 1, autoatención en 2 \\
M4 Rutas / RF-04 & Secuencia y restricciones. & Ruta aprobada a M6. & Planificador de Rutas. & 1 \\
M5 Preparación / RF-05 & Misión y carga confirmada. & Guía nocturna vía \texttt{erp-sync} y ACL en VM-04. & Preparador, Jefa de Bodega. & 1 \\
M6 Reparto / RF-06 & Entrega y POD. & Evento a M7/M8. & Conductor propio, Conductor externo. & 1 \\
M7 Cobranza y rendición / RF-07 & Cobro y descuadre. & ACL al ERP, evento a M10. & Conductores, Gerente de Finanzas. & 1 \\
M8 Devoluciones y envases / RF-08 & Retorno y saldo de envases. & Evento a M2/M7. & Conductores, Jefa de Bodega. & 1 \\
 M9 Calidad y trazabilidad / RF-09  &  Lote, frío, bloqueo y disposición sanitaria autorizada.  &  Estado sanitario a M2/M5, alerta y expediente.  &  Jefa de Calidad.  &  1  \\
M10 Analítica / RF-11 & OTIF operacional y costo de servir. & Consume eventos sin escribir. & Gerentes Comercial, de Finanzas y de Operaciones, Jefa de Calidad. & 1 OTIF, 2 costo \\
M11 Canal moderno / RF-12 & Pedido EDI y excepciones. & Contrato por cadena, M3. & Cliente del canal moderno, Gerente Comercial. & 2 \\
M12 Telemetría / RF-14 & Ruta real, desviación y datos para costo. & GPS a M4/M10. & Planificador de Rutas, Gerente de Operaciones. & 1, costo en 2 \\
\end{tablalafrox}

{\fontsize{9pt}{11pt}\selectfont\raggedright Fuente: elaboración a partir del catálogo de requerimientos y del alcance del capítulo 3. Diseño propuesto: (Subdocumento 4, anexos, diseño de este apartado, pp.~\pageref*{fuente-05-anexo-4-1-d-matriz-de-los-doce-modulos}--\pageref*{fuente-05-anexo-4-1-d-matriz-de-los-doce-modulos-fin}).\par}

\subsection*{Detalle funcional por contexto}
\label{anxD:detalle-contextos}
El catálogo siguiente conserva las capacidades y referencias individuales que sustentan las decisiones resumidas en 4.1.4. Las listas son especificaciones del diseño; no acreditan pruebas ejecutadas ni sustituyen el catálogo completo de requisitos.

\subsubsection*{Módulo M2 Inventario}
\begin{itemize}
  \item \textbf{Stock por sitio (RF-02.02, RT-02.12, Bases Técnicas Transversales, cap. 2, p. 7).} Cada instalación operativa informa sus movimientos al consolidado. Con conexión se consulta disponibilidad actual. Sin ella se muestra la copia descargada, identificada como información con fecha de actualización. La parametrización admite nuevos sitios sin convertir a la casa matriz en una bodega ni asignarle stock operativo por defecto.
\item \textbf{Slotting y conteo ciego.} Asignación de ubicaciones por criterio documentado de rotación, peso y compatibilidad, con conteo cíclico de inventario mediado por terminales HHT con escáner GS1, reduciendo la discrepancia actual del 2,3\% del valor contado a menos del 1\%. El plan mensual de conteo cubre las posiciones que fija Bodega y las elige por rotación, valor y discrepancia histórica. El terminal no muestra la cantidad teórica. M2 compara lo contado con el teórico y registra la diferencia con su causa: si su valor no supera el umbral del recinto, el ajuste se aplica. Si lo supera, la tarea queda pendiente de autorización y se avisa a la jefatura de bodega (RF-02.04, Subdocumento 3, Anexo 3.A).
\item \textbf{FEFO (First Expired, First Out).} La preparación y el despacho respetan el principio de primer vencimiento, con secuenciación térmica para productos refrigerados y congelados. M2 alerta a la jefatura de bodega cuando un lote en stock queda bajo el umbral configurable de vida útil remanente, con la acción recomendada (RF-02.09).
  \item \textbf{Stock disponible.} Con conexión se consulta la disponibilidad mediante APIs. Sin conexión se utiliza la copia autorizada descargada al inicio del turno y almacenada en SQLite/Room. La reserva comercial se confirma en M2 central después del acuse durable de la retención en M2 del sitio. Los pedidos capturados sin señal permanecen a la espera de validación. Los conflictos se resuelven con reglas documentadas de asignación y prioridad, no mediante la sola comparación de marcas de tiempo.

\item \textbf{Cross-docking (RF-02.11).} En cada plataforma, M2 registra la lectura de cada unidad logística al llegar el camión de línea y al cargarla en el camión de reparto, con la hora de cada evento. Una unidad faltante, sobrante o dañada queda como excepción con su causa, de modo que se mide el tiempo real entre consolidación, llegada, desconsolidación y despacho (Subdocumento 3, Anexo 3.A).
  \item \textbf{Reposición de compras (RF-10.01 a RF-10.03).} M2 sugiere la reposición por SKU y proveedor a partir de la venta real, la estacionalidad, las promociones comprometidas y el plazo de cada proveedor. La función de abastecimiento revisa y aprueba la sugerencia con permisos específicos dentro de los perfiles autorizados del CLIENTE. Cada modificación registra responsable y motivo. La sugerencia aprobada se traspasa al módulo de compras del ERP mediante la ACL de INT-06. El ERP conserva la emisión de las órdenes de compra (Subdocumento 3, D-05).
\end{itemize}

\subsubsection*{Módulo M3 Preventa}
\begin{itemize}
  \item \textbf{Toma de pedido offline (RF-03.02).} El preventista toma el pedido sin señal de red. La aplicación mantiene una foto local de datos de lectura (stock, crédito del cliente, precios vigentes, promociones) que se descarga al inicio del turno cuando hay conectividad.
  \item \textbf{Identificador único y deduplicación (RF-03.16).} Cada evento recibe un UUID que se conserva en todos sus reintentos. El servicio de negocio registra su procesamiento y evita ejecutar dos veces la misma operación.
  \item \textbf{Validación de stock y crédito.} La validación definitiva ocurre en la sincronización con la regla de negocio del servidor (reserva de stock, RF-03.03, crédito del cliente, RF-03.08/09). Sin conexión, el pedido queda como a la espera de validación y se resuelve en la sincronización posterior.

\item \textbf{Consulta y control comercial (RF-03.01 a RF-03.10, RF-03.13, RF-03.14 y RF-07.05).} Con conectividad, la aplicación muestra el stock disponible por SKU y el crédito disponible y la deuda vencida del cliente. Sin ella, muestra los de la última sincronización con su antigüedad visible. También muestra el historial de compra con productos y frecuencia, en no más de 2 segundos con conectividad (RNF-03.03). Sugiere productos según historial, frecuencia, productos habituales y promociones vigentes, con los criterios que configura Comercial, y calcula el descuento de la promoción vigente aplicada a una línea. Al seleccionar un cliente con mora de más de 30 días o con su deuda sobre el cupo, la aplicación muestra el bloqueo y solo admite venta al contado o cobro. En los demás casos, si el pedido más la deuda vencida supera el límite de crédito, alerta al preventista antes del envío. Si lo supera en más de 110 \%, bloquea el envío hasta que un supervisor autorizado lo apruebe, con su identidad, fecha y hora en la auditoría. Antes de confirmar, el preventista ve la fecha y la ventana que puede comprometer según RNG-15, el stock reservable, la capacidad de despacho de la ruta y las franjas del cliente (Subdocumento 3, Anexos 3.A y 3.B).
  \item \textbf{Sincronización diferida.} Al recuperar cobertura, el dispositivo envía las escrituras acumuladas por API Gateway. El Gateway aplica los controles de entrada. El servicio de negocio valida, deduplica, procesa y confirma cada evento. El dispositivo conserva los eventos que aún no tienen confirmación.
  \item \textbf{Pedido de autoatención (RT-16.30 y RT-17.01 del caso, Bases Técnicas del caso, cap.~15, p.~27).} M3 recibe también el pedido que el cliente arma en el Portal de Clientes, incluido el armado sin señal en su modo instalable, con la misma deduplicación por UUID y la misma validación de stock y crédito. M3 es el único dueño del pedido en sus tres canales: preventista, autoatención y canal moderno (M11).
\end{itemize}

\subsubsection*{Módulo M4 Rutas}
\begin{itemize}
  \item \textbf{Secuenciación automática.} Optimización determinista de rutas con capacidad y ventanas de tiempo (VRP), con objetivo de ejecución inferior a 20 minutos (Bases Técnicas del caso, cap. 18, p. 34). El planificador puede conocer las restricciones utilizadas y revisar el resultado.
  \item \textbf{Ventanas horarias y capacidad.} Considera capacidad de vehículo, ventanas de entrega del cliente y cadena de frío.
  \item \textbf{Integración con GIS.} Consulta de direcciones, cálculo de ETA y geocercas por API externa, con caché de mapas por zona en dispositivos para degradación a ruta offline con secuencia cargada.
  \item \textbf{Costo de entrega.} Alimenta el cálculo del costo de servir por cliente, uno de los ejes de evaluación del caso.
\end{itemize}

\subsubsection*{Módulo M7 Cobranza y rendición}
\begin{itemize}
  \item \textbf{Rendición digital individual (RF-07.01 a RF-07.03).} M7 calcula para cada camión y turno el monto esperado de rendición con las entregas al contado confirmadas y los cobros registrados en ruta. Cada conductor rinde sus cobros del turno de forma digital contra ese monto, con causales de descuadre clasificadas y trazables. La rendición no se cierra mientras quede un saldo sin una causal tipificada (RF-07.03).
  \item \textbf{Cobranza en ruta (RF-07.06 y RF-07.10).} El conductor registra el cobro y conserva su evidencia para la rendición. En pagos con POS móvil, la captura local de una operación no se presenta como autorización bancaria. El tratamiento sin cobertura depende de las capacidades y condiciones acordadas con la pasarela. Se concilia al reconectar sin generar cargos duplicados. En pagos con tarjeta, la aplicación recibe por Bluetooth el código de autorización del terminal PAX y M7 lo concilia con la entrega en la rendición diaria, sin digitación (RF-07.06). En una entrega al contado, el conductor puede registrar un abono parcial: la aplicación emite un comprobante digital y el saldo insoluto pasa a la cuenta corriente del cliente en el ERP por la ACL (RF-07.10).
  \item \textbf{Pagos informados en el portal (RF-07.08).} Desde la Etapa 2, un pago que el cliente informa en el Portal de Clientes queda «En revisión» en una bandeja donde el cajero de Tesorería verifica el comprobante bancario y autoriza su traspaso estructurado al ERP.
  \item \textbf{Interfaz con ERP (RF-07.09 e INT-06).} La rendición aprobada se publica en la cola FIFO de solicitudes al ERP. Dos procesos del servicio de integración con el ERP en VM-04 de Talca consumen esa cola por conexión saliente y la cola RabbitMQ local, y entregan a la \textbf{capa anticorrupción (ACL)} de la misma VM, la ACL integra con el ERP. Nunca hay escritura directa al ERP. El ERP permanece como única fuente de verdad tributaria y único emisor de DTE.
  \item \textbf{Costo de servir (RF-11).} El módulo alimenta el tablero de BI con el costo real por entrega, incluyendo kilometraje real (telemetría M12), tiempo de servicio y deducciones por devoluciones.
\end{itemize}
\label{fuente-05-anexo-4-1-d-matriz-de-los-doce-modulos-fin}
